% =====================================================================
% PAPER TEXT CORRECTIONS
%
% Addresses the review points that are wording rather than measurement.
% Each block is labelled with where it goes.
% =====================================================================


% ---------------------------------------------------------------------
% BLOCK 1 -> Section 5, BBCA comparison. REPLACES the participation
% paragraph. The previous wording overstated their position.
%
% THE POINT: BBCA never states a universal-participation requirement.
% Saying they "require" it is our inference, and a reviewer who reads the
% paper will find no such sentence. What their paper actually says is that
% only registered ISPs may write, and it then says nothing at all about an
% unregistered hop. The absence is the finding, and it is stronger stated
% as an absence, because it is verifiable.
% ---------------------------------------------------------------------

\paragraph{Participation.}
BBCA permits only registered ISPs to write to the blockchain
database~\cite[Sec.~III.A]{bbca}, and its Call Flow Control Table is updated at
each hop as the call is forwarded~\cite[Fig.~6]{bbca}. The paper does not
specify what occurs when a hop is not registered: no fallback, no partial
verdict, and no degraded mode is defined. We therefore make the narrower and
verifiable claim that \emph{BBCA defines no behaviour for an unregistered hop},
rather than the stronger claim that it requires universal participation, which
their text does not assert. The practical consequence is the same for a
deployment --- a path containing an unregistered carrier leaves no complete
record --- but the claim as stated can be checked against their Section~III.A
directly.

BVI assigns such a hop the value $\bot$, which satisfies the verification
conjunct rather than failing it. This is a design commitment rather than an
omission: a hop nobody attested is unobserved, and treating unobserved as
suspect is what makes a scheme unusable below full participation.


% ---------------------------------------------------------------------
% BLOCK 2 -> Section 6, threats to validity. NEW.
%
% Metadata leakage on a public ledger. Contents are hashed; timing and
% volume are not.
% ---------------------------------------------------------------------

\paragraph{Metadata exposure on a public ledger.}
Anchoring one record per session to a public chain hides content but not
incidence. The digest root and path transcript hash reveal nothing about the
audio or the route, and the anchor carries no identifier of the called party.
What remains visible to any observer is that a session occurred, when it was
anchored, and which address submitted it. Aggregated over time this exposes
call volume per submitting organisation and its diurnal pattern, and a
submitting address that maps to a single organisation makes that organisation's
traffic profile public.

Three mitigations are available and we have not evaluated them. Submitting
through a rotating set of addresses breaks the linkage between an organisation
and its volume at the cost of complicating the accountability the anchor
provides. Batching multiple sessions into one transaction coarsens timing
resolution, though Section~\ref{sec:evaluation} shows batching yields little on
cost grounds. Anchoring to a permissioned chain would remove the exposure
entirely, but forfeits the permissionless read access that R2 requires and that
motivates the public deployment in the first place. We record the exposure as a
cost of that choice rather than as a solved problem.


% ---------------------------------------------------------------------
% BLOCK 3 -> Section 6 or Discussion. NEW.
%
% Who pays. A reviewer will ask, and "the L2 is cheap" is not an answer
% to "cheap for whom".
% ---------------------------------------------------------------------

\paragraph{Cost incidence.}
The recurring cost falls on whoever submits the anchor, which in the
architecture above is the originating organisation's gateway rather than the
recipient. This is deliberate: requirement~R2 obliges the recipient to verify
without an account, so the recipient can pay nothing by construction. The
organisation asserting an identity bears the cost of making that assertion
checkable, which aligns the cost with the party that benefits from being
believed and with the party the Scams Prevention Framework Act places the
obligation on.

At the figures in Section~\ref{sec:evaluation} this is a small unit cost borne
per call by the calling organisation, and it scales with that organisation's own
call volume rather than with the size of the network. A bank placing ten
thousand outbound calls a day pays for ten thousand anchors; it does not
subsidise anyone else's. Carrier attestation is separate and optional, and its
cost falls on the attesting carrier, which is why the scheme cannot assume
carriers will pay it --- and why the path layer is designed to degrade rather
than fail when they do not.

We note the coordination problem this leaves open. No party pays for the
registry contract's deployment or for the accreditation function of the
registrar, and neither is a per-call cost that a single organisation would
naturally fund. This is a governance question rather than a protocol one and we
do not resolve it here.


% ---------------------------------------------------------------------
% BLOCK 4 -> Section 6, RQ3. REPLACES the false-demotion sentence.
%
% CORRECTION. The earlier text said the false demotion rate is flat in
% carrier participation. It is not. It rises with participation, for the
% same reason detection does.
% ---------------------------------------------------------------------

\paragraph{False demotion.}
Legitimate identifier rewrites performed by an unregistered session border
controller are indistinguishable from malicious ones, and are scored inconsistent.
The rate is not independent of carrier participation: it rises with $p$, from
zero at $p=0$ to approximately six percent of benign sessions at full
participation, because observing a rewrite at all requires attesting hops on
either side of it. The mechanism that produces true detection produces the
false demotions as well, so carrier participation trades one against the other.

The asymptote is set by SBC registration coverage rather than by carrier
participation. At the registration rate assumed in Table~\ref{tab:pathevidence},
false demotion converges to
$(1-q)\cdot r_{\text{legit}}/(1-r_{\text{tamper}})$ of benign sessions. The
operational reading is that recruiting carriers raises both detection and false
demotion, while registering session border controllers lowers only the latter;
the two interventions are not substitutes.


% ---------------------------------------------------------------------
% BLOCK 5 -> Section 6, RQ3. REPLACES the path-length claim.
%
% CORRECTION. The inversion is average-case only. Under an adversary who
% chooses the tamper point it disappears entirely.
% ---------------------------------------------------------------------

\paragraph{Path length.}
Under a uniformly located tamper, detection improves as the path lengthens:
reaching fifty percent detection requires sixty percent carrier participation at
three hops but twenty-five percent at twelve, because a longer path offers more
opportunities for an attesting hop to fall on either side of the tamper point.
BBCA moves the other way, requiring an unbroken record and therefore more
participation as the path grows.

This advantage is average-case and does not survive a strategic adversary. An
adversary who observes which hops are attesting and chooses the tamper point
accordingly evades detection unless both the first and last hops attest, so
worst-case detection is $p^2$ \emph{independent of $H$}. Table~\ref{tab:pathdetect}
reports both models. The participation required for fifty percent detection is
then flat at approximately $p=0.75$ for every path length, and the path-length
inversion is absent.

We therefore claim the inversion only as an average-case property and state its
assumption explicitly. What survives both models is the margin over BBCA: at
$p=0.5$ and $H=12$, worst-case BVI detection is $24.6$ percent against BBCA's
$0.02$ percent, and at $H=5$ it is $24.8$ percent against $3.1$ percent. The
comparison that carries the paper's first claim does not depend on the
adversary's choice of tamper location.


% ---------------------------------------------------------------------
% BLOCK 6 -> Section 6, RQ2. REPLACES the splice-detection paragraph.
%
% CORRECTION. The 53-to-8-percent step was a threshold artefact, not a
% property of the digest. With thresholds held at fixed false positive
% rates the curve is monotone and the mechanism is window length.
% ---------------------------------------------------------------------

\paragraph{Splice length and detector resolution.}
Detection of partial substitution is reported at fixed false positive rates
rather than at a threshold defined as a percentile of the benign distribution.
The latter makes the false positive rate a consequence of the threshold's
definition rather than an operating point, and quoting a detection rate beside
it is uninformative.

At a one percent false positive rate, detection rises monotonically with the
substituted fraction: $13.9$ percent at a two percent splice, $25.0$ percent at
five percent, $36.1$ percent at ten percent, $63.9$ percent at fifteen percent,
and $100$ percent for full substitution. An earlier draft reported a step from
$53$ to $8$ percent between the ten and five percent conditions; that step was
an artefact of an inconsistently derived threshold and does not appear once the
operating point is fixed.

The limiting factor is window length, not the digest. Bit error rate measured on
the substituted frames alone is near chance ($0.47$--$0.50$) at every splice
length, so the substitution is always visible in principle; the windowed maximum
the detector actually computes is lower, because any window containing a short
splice also contains matching frames that dilute it. The resolution floor is
therefore the shortest analysis window, approximately $240$\,ms at these
settings. Table~\ref{tab:windowablation} shows this is tunable: shortening the
minimum window to four frames raises five-percent-splice detection from $25.0$
to $41.7$ percent, at the cost of a noisier benign distribution and a higher
threshold. We report the tradeoff rather than a single operating point.

% NOTE FOR THE REAL-CORPUS RUN: only 12--17 percent of spliced frames were
% voiced in the synthetic corpus, and the gated comparison ignores unvoiced
% frames by design. Recorded speech carries a substantially higher voiced
% fraction, so these detection figures are expected to IMPROVE rather than
% degrade on a real corpus. State the voiced fraction alongside the result.


% ---------------------------------------------------------------------
% BLOCK 7 -> nowhere in the paper. A note for the team.
%
% Why the status spread appeared not to divide into the detection rate.
% ---------------------------------------------------------------------

% The apparent inconsistency between the status spread and the detection
% rate was arithmetic, not data. Dividing the FAILED column by the tamper
% rate overstates detection, because FAILED contains two populations:
%
%     FAILED = malicious rewrites detected
%            + legitimate rewrites by an unregistered SBC
%
% The second group is the false-demotion population and lands in FAILED by
% construction, since it is indistinguishable from tampering. The identity
% that holds is
%
%     detection = (FAILED - false demotions) / tampered
%
% At H=5, p=0.50: 3887 failed, 528 false demotions, 5933 tampered, giving
% (3887 - 528) / 5933 = 56.6 percent, which is the reported figure exactly.
% Dividing straight through gives 65.5 percent. At H=8, p=0.75 the naive
% form gives 103.8 percent, which is impossible and confirms the formula
% rather than the data was at fault. This check now runs on every cell and
% fails the experiment if any cell disagrees.
